\section{Related Work}

HIL-RL and intervention learning have developed around the practical need to learn manipulation policies with fewer unsafe or low-value interactions. HIL-SERL-style systems combine human interventions with off-policy reinforcement learning and offline data, while RLPD shows how offline data can be mixed with online updates for efficient learning \cite{luo2025hilserl,ball2023rlpd}. Demonstration-learning infrastructure such as robomimic provides a complementary view of what matters when using human demonstration data for robot manipulation \cite{mandlekar2021robomimic}. These works motivate the learning backbone used here, but they do not by themselves settle how to evaluate a trigger that allocates human attention during training.

Interactive imitation learning and robot-gated correction study related querying problems. ThriftyDAgger treats expert queries as budgeted and gates them by novelty and risk, while adaptive intervention mechanisms can reduce takeover cost in robot-gated imitation learning \cite{hoque2021thrifty,cai2025aim}. Those systems are important baselines for the idea that a robot can decide when to ask, but the current paper studies a different diagnostic question: whether an HIL-RL intervention trigger remains convincing as an attention-allocation policy after cost-matched random baselines and repeated checkpoint estimates.

Model-based and calibration literature informs the diagnostic interpretation. PETS and MBPO motivate uncertainty-aware short rollouts and caution about when learned models should be trusted \cite{chua2018pets,janner2019mbpo}. Calibration work shows that neural-network confidence scores can be miscalibrated, which is directly relevant when a trigger score is used to spend a scarce human budget \cite{guo2017calibration}. Deep RL evaluation studies further warn that small numbers of seeds and point estimates can support unstable conclusions, motivating interval estimates and more careful reporting \cite{henderson2018deeprlmatters,agarwal2021statisticalprecipice}. The present work combines these concerns into an evaluation protocol for intervention-trigger claims that are really claims about human-attention allocation.
